\begin{table}[H]
  \centering
  \small
  \caption{Description textuelle du cas d'utilisation : «~Valider un plan~»}
  \label{tab:cu-valider-plan}
  \PFETableSetup
  \PFETableRowColors
  \PFETableArrayStretch
  \PFETableTabularXBegin
  \begin{tabularx}{\PFETableBlockWidth}{@{}>{\raggedright\arraybackslash}p{3.2cm} >{\raggedright\arraybackslash}X@{}}
    \tableHeaderStyle
    \tableHeaderCell{Champ} & \tableHeaderCell{Contenu} \\

    \textbf{Titre} & Valider un plan \\

    \textbf{Acteurs} &
      1. Utilisateur Corporate.\par
      2. Utilisateur Site (selon le mode de validation). \\

    \textbf{Objectif} &
      Permettre la validation d'un plan en respectant le circuit de validation défini, afin d'assurer la conformité et la cohérence des données avant leur exploitation. \\

    \textbf{Préconditions} &
      1. L'utilisateur est authentifié.\par
      2. Le plan  existe et est en statut «~soumis~».\par
      3. L'utilisateur dispose des droits de validation. \\

    \textbf{Postconditions} &
      1. Le statut du plan est mis à jour.\par
      2. Les informations de validation (date, acteur) sont enregistrées.\par
      3. Le plan est verrouillé après validation complète.\par
      4. Une notification est envoyée aux acteurs concernés. \\

    \textbf{Scénario nominal} &
      1. L'utilisateur accède à la liste des plans soumis.\par
      2. Le système affiche les plans en attente de validation.\par
      3. L'utilisateur sélectionne un plan et lance l'action de validation.\par
      4. Le système vérifie le mode de validation :\par
      \hspace*{1.2em}\textbf{Cas A : Mode validation directe Corporate}\par
      \hspace*{1.2em}4.1 Le système autorise la validation directe.\par
      \hspace*{1.2em}4.2 Le statut du plan passe de Soumis à validé.\par
      \hspace*{1.2em}\textbf{Cas B : Mode validation en deux étapes}\par
      \hspace*{1.2em}4.1 Le manager du site valide le plan.\par
      \hspace*{1.2em}4.2 Le plan est transmis au Corporate.\par
      \hspace*{1.2em}4.3 Le Corporate valide le plan.\par
      \hspace*{1.2em}4.4 Le statut passe à validé.\par
      5. Le système enregistre les informations de validation (acteur, date).\par
      6. Le système verrouille le plan et envoie une notification de validation. \\

    \textbf{Scénario alternatif} &
      \textbf{A. Ordre de validation non respecté :}\par
      A.1 Le Corporate tente de valider avant le manager du site.\par
      A.2 Le système refuse l'opération.\par
      A.3 Un message indique que la validation du site est requise en premier. \\
  \end{tabularx}%
  \PFETableTabularXEnd
\end{table}
